% =====================================================================
%  8. Discussão
%
%  COMPILADO — pontos levantados enquanto as técnicas são percorridas.
%  Nada escrito ainda: cada tópico reúne o achado, sua origem medida e o
%  rótulo onde ele está reportado. A redação vem depois, quando o
%  conjunto estiver fechado.
%
%  A prosa já redigida dos tópicos 1 a 3 está em
%  scratchpad/08-discussao-prosa.tex
% =====================================================================
\section{Discussão}
\label{sec:discussao}


% =====================================================================
%  ESCRITO: Ortografia
% =====================================================================
\subsection{Ortografia}
\label{disc:ortografia}

\indent\indent A dimensão dispõe de um único método, e é ele que devolve o veredito mais forte de todo o
levantamento. A pertinência a $\lexico$ não é estimada nem aproximada: a palavra consta ou não
consta do conjunto, sem gradação e sem margem. É também o único método que, além do
diagnóstico, devolve a correção --- sobre \texttt{natvia}, a Equação~\ref{eq:lex-argmax}
apontou \texttt{nativa} sem ambiguidade, e a subseção~\ref{res:lexico} não registrou falso
positivo sobre a frase correta.

Dois problemas delimitam o que a dimensão alcança, e eles são de naturezas distintas.

O primeiro é de fronteira, e é intrínseco à definição. Em \emph{A três dias} o método não
produziu achado, porque \texttt{a} é palavra do português e pertence a $\lexico$. A Equação~\ref{eq:lex-deteccao} avalia uma palavra por
vez, isoladamente, e \texttt{a}, tomada isoladamente, não apresenta desvio algum.

O segundo problema é de procedência, e afeta o que o método devolve quando acerta. O léxico é
uma lista de frequências, e não um dicionário: o conjunto de candidatos herda o que o corpus
de origem contém. Dos quatro candidatos a uma edição de \texttt{natvia} na
Tabela~\ref{tab:candidatos}, apenas \texttt{nativa} é palavra corrente do português;
\texttt{naevia} e \texttt{natia} não constam de dicionário normativo, e \texttt{latvia} é
forma estrangeira. São candidatos por terem sido atestados no material de origem, e não por
serem palavras do idioma. O desempate segue a mesma lógica, uma vez que a decisão recai
integralmente sobre $f$: \texttt{nativa} vence com 912 ocorrências, então, a qualidade 
do resultado esta diretamente ligada ao corpus utilizado.



% =====================================================================
%  ESCRITO: Gramática
% =====================================================================
\subsection{Gramática}
\label{disc:gramatica}

\indent\indent A dimensão reúne quatro métodos, e eles se dividem em dois grupos pelo tipo de resposta. Os
traços morfossintáticos e as regras escritas à mão devolvem um veredito localizado; o SLOR e o
PLL devolvem um número.

Os dois primeiros acertam onde têm cobertura, sendo assim, dependem das regras bem descritas.
 As regras detectaram a troca em \emph{A três dias} e são o único método do
levantamento a devolver, junto com o achado, o enunciado da norma em português; mas sobre
\emph{os incêndios atingiu} nenhuma das 2.270 regras casou, e $D(\sent) = \varnothing$ é
indistinguível da saída de uma frase correta. É a Equação~\ref{eq:regras-invisivel} verificada
em caso concreto: a cobertura de $R$ é finita e de elaboração manual, e o português dispõe de
249 grupos de regras. Os traços, por sua vez, localizaram o desvio de
\emph{uma áreas} com precisão que nenhum outro método alcança --- a palavra, o termo com que
diverge, o traço violado e o arco que os liga ---, mas deixaram passar \emph{a área inteiro}:
o adjetivo posposto está no masculino e ainda assim $E(\sent) = \varnothing$, porque o
anotador atribuiu $\texttt{Gender} = \texttt{Fem}$ a \texttt{inteiro} e a comparação devolveu
\emph{iguais}. A fórmula está correta e o insumo não está, e é por isso que a
Equação~\ref{eq:tracos-E} compara $\tracosp{}{}$, e não $\tracos{}{}$.

Os dois métodos de escore ordenam, mas não decidem. O PLL ordenou as quatro frases na direção
certa, com degraus amplos, e ainda localizou o desvio, uma vez que $5{,}18$ dos $6{,}00$ nats
de diferença se concentram na posição do verbo. Essa localização, porém, é atribuição do
modelo e não fato da frase: a discordância é \emph{entre} sujeito e verbo, e mascaradas as
formas \texttt{os} e \texttt{incêndios} tornam-se mais prováveis que suas contrapartes no
singular. Nenhum dos dois métodos dispõe de limiar, de modo que os valores são interpretáveis
entre si e nenhum isoladamente.

No SLOR soma-se um problema do estimador. O par de ordem foi separado na direção prevista, mas
a Tabela~\ref{tab:slor} mostra sobre o que a separação repousa: nove das onze parcelas da
frase embaralhada estão no piso $\ln(1-\lambda) = -1{,}204$, e a única contribuição positiva
vem de \texttt{de o}, adjacência entre preposição e artigo que sobreviveu ao embaralhamento
por acidente; do outro lado, 43\% do numerador da frase correta provém de um único bigrama
observado duas vezes. O escore não distingue frase mal construída de frase cujo vocabulário o
corpus de contagens não cobre, e assim nenhum dos quatro métodos decide sozinho sobre
concordância, dois são dependente de regras que podem falhar em casos de cenários 
fora de sua cobertura e dois não dispõem de limiar.



% =====================================================================
%  ESCRITO: Estrutura
% =====================================================================
\subsection{Estrutura}
\label{disc:estrutura}

\indent\indent Os dois métodos da dimensão ocupam posições opostas quanto ao que se pode extrair deles. O
teste da raiz devolve veredito binário e localizado ao custo de inspecionar um único nó; o
comprimento de dependência e a não-projetividade devolvem descritores que não respondem à
pergunta da correção.

O teste da raiz registra uma distinção que nenhum outro método do levantamento alcança. Entre
\emph{o incêndio atingiu} e \emph{o incêndio que atingiu} não há diferença de vocabulário ---
o mesmo verbo está presente nas duas ---, e o que separa oração de fragmento é a posição que
ele ocupa na árvore. Das quatro frases a que foi submetido, o método acertou três, aprovando a
frase de referência pelo primeiro termo da disjunção e a predicação nominal pelo segundo, e
reprovando o fragmento por ambos serem falsos.

O erro é de espécie mais simples que os da dimensão anterior. Em \emph{Atingir uma área de
mata nativa em três dias} a raiz é \texttt{Atingir}, também VERB, e o infinitivo isolado é
aprovado como oração. O traço que separaria os dois casos ---
$\texttt{VerbForm} = \texttt{Inf}$ --- está na árvore, corretamente atribuído, e a
Equação~\ref{eq:raiz-pred} simplesmente não o consulta. Diferentemente do que ocorreu com os
traços morfossintáticos, aqui o insumo está correto e a falha é da definição, portanto
corrigível dentro dela.

Uma ressalva recai sobre a árvore que o método lê. A notação $\raiz(\sent)$ pressupõe
unicidade, e o analisador não a garante: sobre o fragmento devolveu uma floresta, com
\texttt{dias} e \texttt{incêndio} ambos etiquetados \texttt{ROOT}. A implementação retém o
primeiro, e no caso o veredito não depende da escolha, pois os dois candidatos são NOUN --- mas
o acerto se deu sobre uma entrada que a definição não previa.

O segundo método não emite veredito algum, e o corpus deixa isso explícito: as duas frases a
que foi aplicado são corretas e diferem apenas na posição em que as relativas foram
encaixadas. O comprimento médio as separa com margem ampla, de $1{,}42$ para $3{,}58$, com o
maior arco passando de 2 para 9, o que mensura a carga de processamento imposta pela ordem
escolhida. Já a não-projetividade devolveu $X = 0$ nas duas: as desigualdades da
Equação~\ref{eq:dep-cruzam} são estritas, e o encaixe de uma relativa no interior de outra
produz arcos aninhados, não sobrepostos. A medida está correta e sobre este corpus não informa
nada, de modo que da dimensão inteira apenas o teste da raiz é aproveitável para decidir sobre
uma frase.



% =====================================================================
%  ESCRITO: Semântica
% =====================================================================
\subsection{Semântica}
\label{disc:semantica}

\indent\indent A dimensão reúne três métodos, e apenas um deles produziu decisão utilizável.

O verbo mascarado separou o verbo anômalo dos adequados por quase 14 nats, e o zero de sua
escala tem significado, uma vez que \texttt{destruiu} é a primeira escolha do modelo e recebe
margem exatamente nula. Mais relevante que a separação, porém, é a condição que a torna
legítima: a lista de formas esperadas é idêntica nas três frases, porque mascarar a raiz
suprime justamente a única palavra em que elas diferem. O instrumento de medida permanece fixo
enquanto o objeto varia --- propriedade que nenhum outro método do levantamento possui, dado
que no SLOR e na surpresa trocar a palavra altera também as contagens que servem de referência.

Duas ressalvas acompanham o método. A resposta \texttt{n/a} sobre o fragmento e a predicação
nominal é comportamento previsto, e não falha, mas transfere a decisão ao analisador
sintático: as duas frases foram recusadas porque a árvore predita as classificou assim. E
$\mathcal{V}$ não está restrito a formas verbais, como mostra a preposição \texttt{em} entre
as cinco esperadas; neste corpus a folga é inconsequente, pois $\hat{v}$ foi \texttt{destruiu}
nas três frases, mas nada na Equação~\ref{eq:verbo-margem} impede que o melhor preenchimento
seja palavra de outra classe.

A surpresa por token acertou uma das três frases e errou duas, e os três picos têm a mesma
origem. Todos recaem sobre posições de bigrama não observado, onde vale
$\surp_{\mathrm{uni}} + 1{,}204$: na frase de referência, que não contém desvio, o pico caiu
sobre \texttt{atingiu} apenas porque \texttt{incêndio atingiu} não consta do corpus de
contagens. Em \emph{Focos recorrentes apresentaram persistência anômala} o caso é extremo,
pois os cinco bigramas são não observados, todo o perfil é unigrama deslocado de uma constante,
e o pico é a palavra mais rara no valor de teto do estimador. O $\arg\max$ seleciona
raridade, e o acerto sobre \texttt{cantou} é contingência de o desvio ser também uma palavra
rara.

Os embeddings falham por motivo distinto. O bloco de pares sem relação comporta-se como
esperado, situando-se junto a zero, mas o bloco superior não separa o que deveria: o par de
maior proximidade de toda a Tabela~\ref{tab:embeddings} é um par de opostos,
\texttt{seco}~×~\texttt{úmido} com $+0{,}838$, acima de todos os sinônimos e do próprio
\texttt{incêndio}~×~\texttt{fogo}. Não é defeito do modelo, e sim a hipótese distribucional
operando como definida, pois os dois termos ocorrem nos mesmos contextos. Some-se que a
Equação~\ref{eq:vetor-frase} é invariante à ordem --- verificada em $1{,}000000$ sobre o par
de permutação --- e que \texttt{natvia}, fora da tabela de vetores, não recebe medida alguma.

Das três, portanto, apenas o verbo mascarado sobrevive à caracterização, e as duas exclusões se
dão por razões opostas: a surpresa mede raridade em lugar de desvio, e o cosseno mede campo
temático em lugar de sentido.



% =====================================================================
%  ESCRITO: A pipeline (introdução e seleção)
% =====================================================================
\subsection{A pipeline}
\label{disc:pipeline}

\indent\indent Tendo em vista os resultados obtidos, 
observa-se uma abordagem por partição: cada um alcança um tipo 
de desvio e falha em outros. Realizando uma análise segmentada, 
podemos priorizar técnicas em que as falhas não coincidem. 
--- \emph{A três dias} escapa ao léxico e é detectada pelas regras; \emph{os
incêndios atingiu} escapa às regras e é detectada pelo PLL. Encadear os métodos é a resposta
que decorre dessa observação, e o que se apresenta a seguir é uma sugestão de cadeia, não um
resultado do levantamento.

Propõem-se então quatro etapas, uma por dimensão, cada qual ocupada pelo método que a discussão
anterior identificou como o único da sua dimensão a produzir veredito localizado sem parâmetro
a calibrar:

\begin{center}
\small
ortografia $\;\rightarrow\;$ estrutura $\;\rightarrow\;$ gramática $\;\rightarrow\;$ semântica \\[2pt]
\emph{léxico} \qquad \emph{teste da raiz} \qquad \emph{traços} \qquad \emph{verbo mascarado}
\end{center}

A primeira etapa é o léxico porque é o único veredito do levantamento que não depende de
estimativa: a palavra consta ou não consta de $\lexico$. É também o único método que devolve a
correção junto com o diagnóstico. A ressalva da subseção~\ref{disc:ortografia} permanece --- a
qualidade do que ele devolve está ligada ao corpus que originou a lista ---, mas ela afeta a
correção sugerida, não a detecção.

A segunda é o teste da raiz porque, das duas técnicas estruturais, é a única que emite
veredito: o comprimento de dependência mede carga de processamento sobre frases corretas, e a
não-projetividade não separou nada. Pesa a favor dele que sua única falha, a aprovação do
infinitivo isolado, seja corrigível dentro da própria definição, bastando que a
Equação~\ref{eq:raiz-pred} passe a exigir $\texttt{VerbForm} = \texttt{Fin}$ na raiz. É a
correção que a cadeia pressupõe.

A terceira são os traços morfossintáticos, e não as regras escritas à mão, embora as duas
detectem desvios reais. A razão é a forma da cobertura: $T(r)$ é uma condição sobre classes de
arcos e vale para toda ocorrência da classe, ao passo que $R$ é um conjunto finito de padrões
individuais, cuja ausência de item torna o desvio invisível. Os traços também localizam com a
precisão mais fina do levantamento. As regras permanecem, contudo, a extensão mais imediata da
cadeia, uma vez que detectam a escolha lexical, que é o desvio que nenhuma das quatro etapas
alcança.

A quarta é o verbo mascarado por ser o único método semântico que sobreviveu à caracterização,
e porque é o único de todo o levantamento cujo instrumento de medida permanece fixo enquanto o
objeto varia.

A ordem das quatro não é escolha de projeto. A ortografia vem primeiro porque uma palavra fora
do léxico não permanece contida nessa etapa: ela deforma a árvore de dependências que as três
etapas seguintes leem. A estrutura precede a gramática porque a existência de predicado é
condição para que os arcos de concordância estejam corretamente atribuídos, e o teste da raiz a
verifica a custo desprezível. E a semântica encerra a cadeia porque o verbo mascarado devolve
\texttt{n/a} quando a raiz não é verbo, de modo que depende da verificação estrutural. Cabe
registrar que essa dependência não é estrita: uma predicação nominal passa pela segunda etapa e
ainda assim recebe \texttt{n/a} na quarta, pois ter predicado e ter verbo na raiz não são a
mesma condição.


Aplicada a cadeia, cada etapa devolve um achado ou nada, e a frase é reprovada quando alguma
delas produz achado. As quatro são executadas sempre, ainda que a primeira já reprove: a
decisão de não bloquear é o que torna o laudo legível, pois expõe todos os desvios de uma vez
em lugar de um por passagem. A Tabela~\ref{tab:laudo} apresenta o laudo de duas frases do
corpus.

\begin{table}[H]
\centering
\dimtable{
\begin{tabular}{@{}llll@{}}
\toprule
& \textbf{Etapa} & & \textbf{Achado} \\
\midrule
\multicolumn{4}{@{}l}{\emph{Há três dias o incêndio atingiu uma áreas de mata nativa.}} \\
\midrule
& ortografia & ok & --- \\
& estrutura  & ok & predicado presente \\
& gramática  & \textbf{!!} & Number: \texttt{uma}(Sing) × \texttt{áreas}(Plur) [det] \\
& semântica  & ok & \texttt{atingiu}, margem 0,65 \\
\midrule
\midrule
\multicolumn{4}{@{}l}{\emph{Há tres dias o incêndio atingiu uma áreas de mata natvia.}} \\
\midrule
& ortografia & \textbf{!!} & \texttt{natvia} $\to$ \texttt{nativa} \\
& estrutura  & ok & predicado presente \quad$^\dagger$ \\
& gramática  & \textbf{!!} & Number: \texttt{uma}(Sing) × \texttt{áreas}(Plur) [det] \quad$^\dagger$ \\
& semântica  & ok & \texttt{atingiu}, margem 0,95 \quad$^\dagger$ \\
\bottomrule
\end{tabular}}
\caption{Laudo da cadeia sobre duas frases. $^\dagger$ etapa marcada como \emph{árvore sob
suspeita}, por ler uma árvore construída sobre palavra fora do léxico.}
\label{tab:laudo}
\end{table}

A primeira frase contém um único desvio, e a cadeia o localiza na etapa que lhe corresponde. A
segunda acrescenta uma palavra fora do léxico, e é ela que expõe o custo do encadeamento: a
forma \texttt{natvia} não permanece contida na primeira etapa, pois deforma a árvore de
dependências que as três seguintes leem. Daí a marcação \emph{árvore sob suspeita}, que não
indica incorreção, mas confiabilidade menor --- e que é a resposta prática ao limite enunciado
na subseção~\ref{disc:gramatica}, o de que uma cadeia de etapas que leiam a mesma árvore
predita herda os erros dessa predição.

No caso examinado a deformação não atingiu o arco relevante, e o achado de concordância
permanece legítimo nas duas frases. Trata-se, contudo, de contingência posicional, e não de
garantia: a marcação existe justamente porque a cadeia não tem como distinguir os dois casos.

Cabe observar, por fim, que a quarta etapa não produz achado em nenhuma das duas frases, com
margens de $0{,}65$ e $0{,}95$. O ganho dela não se manifesta neste par, cujos desvios são de
ortografia e de concordância; ela se destina à anomalia de seleção verbal, para a qual a margem
medida na subseção~\ref{res:verbo} alcança $13{,}73$.



% =====================================================================
%  TÓPICO 1 — [ESCRITO em \subsection{Ortografia}]
%  origem: res:lexico, tab:candidatos
% ---------------------------------------------------------------------
%  - O léxico é lista de frequências, não dicionário. O conjunto de
%    candidatos herda o que o corpus contém: dos quatro candidatos de
%    'natvia', só 'nativa' é palavra corrente do português — 'naevia' e
%    'natia' não constam de dicionário normativo, 'latvia' é estrangeira.
%  - O desempate reproduz a distribuição do corpus: 'nativa' vence com
%    f = 912, margem do material de origem e não propriedade linguística.
%  - Trocar a lista altera quem entra em C(w) E quem vence o argmax, sem
%    que uma linha da definição mude.
%  - Consequência: a qualidade da correção é herdada, não construída. O
%    veredito de pertinência a D — o único exato entre os métodos — é
%    exato apenas relativamente ao conjunto que se adotou.


% =====================================================================
%  TÓPICO 2 — Três falhas, três causas
%  origem: res:lexico, res:regras, e o caso 'a área inteiro' (célula 9)
% ---------------------------------------------------------------------
%  Cada método deixou passar um desvio, e as causas estão em níveis
%  diferentes. É isso que justifica encadear em vez de aperfeiçoar um.
%
%  (a) DEFINIÇÃO — 'A três dias', no léxico.
%      'a' pertence a D; eq:lex-deteccao se esgota nesse fato. Não há
%      defeito a corrigir: predicado que avalia uma palavra isolada não
%      distingue forma existente de forma adequada.
%
%  (b) COBERTURA DO RECURSO — 'os incêndios atingiu', nas regras.
%      Nenhuma das 2.270 regras casou; D(s) = vazio. eq:regras-invisivel
%      previa exatamente isso. Correção disponível: escrever mais uma
%      regra — trabalho manual que não generaliza.
%
%  (c) INSUMO — 'a área inteiro', nos traços.  <-- a mais difícil
%      Adjetivo posposto no masculino, núcleo no feminino, e ainda assim
%      E(s) = vazio. O arco amod foi admitido por T(r), a comparação de
%      Gender foi efetuada e devolveu 'iguais': o anotador atribuíra
%      Gender=Fem a 'inteiro'. eq:tracos-E compara mu-chapéu, não mu.
%      A formulação está correta; nenhum ajuste dela recupera o caso.
%      É a promessa que met:tracos deixa em aberto e que aponta para cá.


% =====================================================================
%  TÓPICO 3 — O que as falhas justificam, e o que elas não justificam
%  origem: res:lexico + res:regras (par medido de complementaridade)
% ---------------------------------------------------------------------
%  - Nenhuma das três falhas se corrige por dentro do método que a
%    produziu, e elas não coincidem: 'A três dias' escapa ao léxico e é
%    detectada pelas regras. É o argumento direto para o encadeamento —
%    não porque um método seja melhor, mas porque operam sobre objetos
%    distintos e erram em lugares distintos.
%  - LIMITE, a enunciar junto: a falha (c) é a que o encadeamento NÃO
%    resolve. Etapas que leiam a mesma árvore predita herdam o erro dessa
%    predição. O encadeamento resolve complementaridade de cobertura, não
%    qualidade do insumo.
%  - Daí a necessidade de uma etapa registrar quando há motivo para
%    suspeitar da árvore, em vez de tratar todo veredito como igualmente
%    confiável. (Prepara a marcação 'árvore sob suspeita' da pipeline.)


% =====================================================================
%  TÓPICO 4 — O SLOR não ordena boa formação
%  origem: res:slor, tab:slor
% ---------------------------------------------------------------------
%  - PAR MÍNIMO SEM SEPARAÇÃO: 'Os satélites registraram/registrou 47
%    focos' recebem o MESMO escore, -0,072, até a 3a casa. Nenhum dos dois
%    bigramas foi observado; ambas as posições recuam para a unigrama.
%  - FRASE CORRETA NO PISO: 'Focos recorrentes apresentaram persistência
%    anômala' = -1,204 = ln(1-lambda), exatamente eq:piso-slor. Nenhum de
%    seus 5 bigramas observado. Mede ausência de vocabulário no corpus de
%    contagens, não construção.
%  - ORDENAÇÃO INVERTIDA: a frase correta pontua ABAIXO da embaralhada
%    (-1,000) e ABAIXO da frase sem sentido de Chomsky (-0,939).
%  - TESE: a raridade que o método foi construído para descontar reaparece
%    pelo piso do estimador. Liga ao TÓPICO 1 — de novo o corpus de
%    contagens determinando o veredito, agora num método de escore.
%  - Consequência para a cadeia: escore global sem termo de comparação não
%    decide sobre frase isolada. Reforça o veredito por vacuidade de
%    conjunto em vez de escore agregado.


% =====================================================================
%  A ACRESCENTAR conforme as técnicas forem percorridas:
%  PLL, raiz, dependência, surpresa, verbo, embeddings
% =====================================================================
